\section{Introduction}
Procedural Content Generation (PCG) has been actively researched in computer graphics for years. For developers generating virtual content can become time-intensive, and PCG has been used to reduce the workload of creating virtual terrains. One important aspect is the trade-off between the control the user has over the generated terrain against the amount of effort required to create the terrain. For this research, we will be designing a system to generate complete landscapes which will be rendered from ground level. Besides combining multiple features, this research will aim to give a qualitative evaluation of the generated terrain with rules that can be specified by the designer. This system will then be used to create landscapes as part of the Growing Roots research project.
\todo{Show how this system will be different than other implementations.}